%=============================================================
% APPENDIX
%=============================================================

\appendix

\section{Supplementary Information of Zero-Trust Memory Game}
\label{app:problem-def}

\subsection{Full Rationale for the Zero-Trust Formulation}
\label{app:zt_rationale}

The zero-trust formulation is motivated by three converging lines of evidence.

\textbf{Misalignment in LLM-based agents.}
Large language model agents have been shown to exhibit sycophantic
drift~\citep{fanous2025syceval,sharma2023towards}, wherein a model progressively
adjusts its stated beliefs toward those of a conversational partner regardless
of evidential support. In the context of shared memory, this implies that a
contributor agent may craft $\Delta_t$ not to reflect its best epistemic state
but to conform to or reinforce a perceived consensus already encoded in
$\mathcal{M}_t$. Analogously, an auditor agent subject to sycophancy will
approve $\Delta_t$ because it matches existing memory patterns rather than
because it is true. Both behaviors are misaligned with collective memory
integrity yet neither constitutes classic Byzantine behavior, i.e., the agent is not
``faulty'' in the sense of crashing or sending random messages; it is behaving
systematically in a way that is undetectable by fault-tolerance mechanisms.

\textbf{Limitations of Byzantine fault tolerance.}
Byzantine fault tolerance (BFT)~\citep{castro1999practical,lamport2019byzantine}
assumes that up to $f < n/3$ agents may behave arbitrarily, while the rest are
honest. This framework is ill-suited to multi-agent LLM systems for two reasons.
First, the fraction of misaligned agents is not known in advance and may exceed
the BFT threshold, particularly when agents share the same base model and
therefore share systematic biases~\citep{turpin2023language}. Second, BFT
assumes honest agents are identifiable by their consistent and correct behavior,
whereas in LLM systems the same agent may behave honestly in one context and
misalign in another depending on the prompt, the memory state, or the identities
of other agents~\citep{perez2022red}. The zero-trust assumption removes the
dependency on a trusted majority entirely.

\textbf{Social choice and manipulation in multi-agent deliberation.}
Multi-agent debate frameworks~\citep{du2024improving, liang2024encouraging,
chan2023chateval, xiong2023examining, wang2024rethinking}
aggregate reasoning by having agents critique and revise each other's outputs.
Results in social choice theory~\citep{arrow2012social, gibbard1973manipulation,
satterthwaite1975strategy} establish that any aggregation rule over rational agents
is susceptible to intentional manipulation by at least one participant, where there is
no mechanism that is simultaneously strategy-proof, Pareto-efficient, and
non-dictatorial (Arrow's impossibility theorem \citep{geanakoplos2005three}). Applied to memory updating,
this implies that for any commitment rule $\mathcal{C}$, there exists a profile
of agent strategies under which the committed memory update is not the
collectively optimal one. Zero-trust is therefore not a pessimistic special case
but the general condition under which shared-memory deliberation operates.

\subsection{Rationale for Utility Structures}
\label{app:utility_rationale}

\textbf{Why the contributor utility takes a product form?}
The contributor's utility $u_c = \Pr[\mathcal{C}=1] \cdot \mathbb{I} - D$
is designed to reflect a strategic reality: committing an uninfluential
delta has no value, and committing a high-influence delta that is immediately
detected has negative value. The product $\Pr[\mathcal{C}=1] \cdot
\mathbb{I}$ captures this joint requirement, a contributor must
simultaneously achieve commitment \emph{and} influence to gain positive
utility. Removing the scalar weights $\mu_c$ and $\nu_c$ from the earlier
formulation is justified by noting that their ratio is not identifiable
from observed behavior: only the tradeoff between the product term and $D$
matters for equilibrium characterization, so we normalize without loss of
generality.

\textbf{Why auditor utility measures pre-commit alignment?}
The original formulation $\phi_i(\mathcal{M}_{t+1})$ conditions on the
post-commit memory state, which is circular: the auditor's vote determines
whether the commit occurs, so its utility depends on its own action in a
way that conflates strategic incentives with outcome evaluation. Instead,
$\texttt{Align}(\Delta_t, \mathcal{M}_t)$ is a pre-commit signal that the
auditor observes \emph{before} voting, making the utility well-defined as
an expected-payoff function over the auditor's strategy. This is consistent
with standard extensive-form game formulations where payoffs are evaluated
at terminal nodes given the full history, but the relevant information for
strategy selection is what is observable before the action is taken.


\textbf{Online estimation of $\texttt{Align}(\cdot)$ and $c_i$.}
The two remaining quantities in the auditor utility are estimated without
human pre-specification.

\textit{Alignment signal.} $\texttt{Align}(\Delta_t, \mathcal{M}_t)$ is
computed by $\Phi$ as a normalized version of the architecture-specific
calibration score:
\begin{equation}
    \texttt{Align}(\Delta_t, \mathcal{M}_t) =
    2\rho(\Delta_t, \mathcal{M}_t) - 1 \in [-1, +1]
\end{equation}
where $\rho \in \{\rho^{\text{emb}}, \rho^{\text{graph}}\}$ is the
calibration score defined in Section~\ref{sec:calibration}. This
grounding means $\texttt{Align}$ is entirely determined by the structural
properties of $\Delta_t$ relative to $\mathcal{M}_t$, with no free
parameters.

\textit{Effort cost.} The effort cost $c_i$ is estimated as the
empirical mean token consumption (or wall-clock time) of $a_i$'s
scrutiny actions across the most recent $W$ rounds:
\begin{equation}
    \hat{c}_i^{(t)} = \frac{1}{W} \sum_{t'=t-W}^{t-1}
    \text{cost}(a_i,\ t')
\end{equation}
where $\text{cost}(a_i, t')$ is the measured computational cost of
$a_i$'s participation in round $t'$. This is observable without any
assumptions about agent internals and requires no human input. The
window size $W$ is set to the minimum value satisfying a standard
variance criterion: $W^* = \min\{W : \text{Var}(\hat{c}_i^{(t)}) \leq
\varepsilon_c\}$, where $\varepsilon_c$ is the acceptable estimation
error, making even this meta-parameter data-driven.

\subsection{Equilibrium Failure Modes Under Zero-Trust}
\label{app:failure_modes}

We formally establish three canonical failure modes that emerge as Nash
equilibria of $\Gamma$ under Assumption~\ref{assm:zerotrust}. These
propositions collectively motivate why memory integrity cannot be achieved
by agent-level trust alone, and provide the formal basis for the
calibration mechanism of Section~\ref{sec:calibration}.

\begin{proposition}[Auditor Free-Riding]
\label{prop:freeriding}
Under a majority commitment rule $\mathcal{C}$ with $|\mathcal{A}_t| = n$
auditors, if each auditor's effort cost satisfies $c_i > 0$ and the
alignment signal $\emph{Align}(\Delta_t, \mathcal{M}_t)$ is independent
of any individual auditor's vote, then the unique dominant-strategy
equilibrium is one of \emph{rubber-stamp approval}: each $a_i \in
\mathcal{A}_t$ emits $v_i = 1$ without scrutiny, regardless of $\Delta_t$.
\end{proposition}

\begin{proof}
Fix any auditor $a_i \in \mathcal{A}_t$ and let $\mathbf{v}_{-i}$ denote
the votes of all other auditors. Under majority rule, $a_i$'s vote is
pivotal (i.e., changes the commitment outcome) only if exactly
$\lfloor n/2 \rfloor$ of the remaining $n-1$ auditors vote to approve.
Let $p_{\text{piv}}$ denote this pivotality probability. For any symmetric
mixed strategy profile over the other auditors with approval probability
$q \in [0,1]$:
\begin{equation}
    p_{\text{piv}} = \binom{n-1}{\lfloor n/2 \rfloor}
    q^{\lfloor n/2 \rfloor}(1-q)^{\lceil n/2 \rceil - 1}
    = O(n^{-1/2})
\end{equation}
by Stirling's approximation. The marginal integrity gain from scrutiny is
therefore $p_{\text{piv}} \cdot \Delta\texttt{Align} \leq p_{\text{piv}}
\cdot 2$, since $\texttt{Align} \in [-1,+1]$. The expected net utility gain
from scrutiny over rubber-stamping is:
\begin{equation}
    \Delta u_i = p_{\text{piv}} \cdot \Delta\texttt{Align} - c_i
    \leq 2 p_{\text{piv}} - c_i = O(n^{-1/2}) - c_i
\end{equation}
Since $c_i > 0$ is fixed and $p_{\text{piv}} \to 0$ as $n \to \infty$,
there exists $n^* = O(c_i^{-2})$ such that $\Delta u_i < 0$ for all
$n > n^*$. For any $n \geq 2$, $v_i = 1$ without scrutiny weakly
dominates careful scrutiny because the effort cost $c_i$ is incurred
regardless of whether $a_i$ is pivotal, while the integrity payoff
from scrutiny requires pivotality. Therefore, $v_i = 1$ is the weakly
dominant strategy for all $a_i$, and the unique dominant-strategy
equilibrium is universal approval without scrutiny.
\end{proof}

\begin{proposition}[Collusion Equilibrium]
\label{prop:collusion}
Suppose a coalition $\mathcal{K} \subseteq \mathcal{A}_t$ of auditors
shares a positive collusion benefit $\lambda_i > 0$ from approving
$\Delta_t$ regardless of its content, and $|\mathcal{K}| \geq
\lceil |\mathcal{A}_t| / 2 \rceil$. Then there exists a Nash equilibrium
in which every $a_i \in \mathcal{K}$ approves $\Delta_t$ unconditionally,
the commitment rule $\mathcal{C}$ is satisfied, and the committed delta
may be arbitrarily inconsistent with $\mathcal{M}_t$.
\end{proposition}

\begin{proof}
We construct the equilibrium strategy profile explicitly. Let each
$a_i \in \mathcal{K}$ play the pure strategy $v_i = 1$ unconditionally,
and let agents outside $\mathcal{K}$ play any strategy. We verify that
no agent in $\mathcal{K}$ has a profitable deviation.

\textit{Payoff under collusion.} For $a_i \in \mathcal{K}$, playing
$v_i = 1$ yields:
\begin{equation}
    u_i^{\text{collude}} = \texttt{Align}(\Delta_t, \mathcal{M}_t)
    + \lambda_i
\end{equation}
where the effort cost term is zero since no scrutiny is performed.

\textit{Payoff under deviation.} If $a_i$ deviates to $v_i = 0$
(rejection), the collusion benefit $\lambda_i$ is forfeited. Since
$|\mathcal{K}| \geq \lceil |\mathcal{A}_t|/2 \rceil$, the majority
threshold is still satisfied by $\mathcal{K} \setminus \{a_i\}$
whenever $|\mathcal{K}| > \lceil |\mathcal{A}_t|/2 \rceil$. In this
case the deviation does not change the commitment outcome and the
deviating auditor loses $\lambda_i > 0$:
\begin{equation}
    u_i^{\text{deviate}} = \texttt{Align}(\Delta_t, \mathcal{M}_t)
    - c_i < u_i^{\text{collude}}
\end{equation}
When $|\mathcal{K}| = \lceil |\mathcal{A}_t|/2 \rceil$ exactly,
deviation by $a_i$ does change the outcome to rejection, but
$a_i$'s payoff then becomes $\texttt{Align}(\Delta_t, \mathcal{M}_t)
- c_i - \lambda_i < u_i^{\text{collude}}$ since $\lambda_i > 0$ and
$c_i > 0$. In both cases deviation is strictly dominated.

\textit{Commitment outcome.} Since $|\mathcal{K}|$ meets the majority
threshold, $\mathcal{C}(\{v_i\}_{i \in \mathcal{A}_t}) = 1$
regardless of the content of $\Delta_t$. In particular, $\Delta_t$
may satisfy $\texttt{Align}(\Delta_t, \mathcal{M}_t) = -1$ (maximal
inconsistency) and still be committed. No agent in $\mathcal{K}$ has
incentive to deviate, so this constitutes a Nash equilibrium.
\end{proof}

\begin{proposition}[Adversarial Rejection]
\label{prop:adversarial_rejection}
Suppose a blocking coalition $\mathcal{B} \subseteq \mathcal{A}_t$
satisfies $|\mathcal{B}| > |\mathcal{A}_t|/2$, and each $a_i \in
\mathcal{B}$ has utility strictly decreasing in the contributor $a_c$'s
influence on $\mathcal{M}$, i.e., $u_i$ is decreasing in
$\mathbb{I}(\Delta_t, \mathcal{M}_t)$ for any $\Delta_t$ from $a_c$.
Then there exists a Nash equilibrium in which all $\Delta_t$ from $a_c$
are rejected regardless of their consistency with $\mathcal{M}_t$.
\end{proposition}

\begin{proof}
Let each $a_i \in \mathcal{B}$ play the pure strategy $v_i = 0$
unconditionally for any $\Delta_t$ proposed by $a_c$.

\textit{Payoff under adversarial rejection.} For $a_i \in \mathcal{B}$,
rejecting $\Delta_t$ prevents commitment and therefore prevents $a_c$
from gaining influence over $\mathcal{M}$. Since $u_i$ is decreasing
in $\mathbb{I}(\Delta_t, \mathcal{M}_t)$, blocking the commit yields:
\begin{equation}
    u_i^{\text{reject}} = \texttt{Align}(\Delta_t, \mathcal{M}_t)
    \cdot \mathbf{1}[\mathcal{C} = 1] - c_i \cdot
    \mathbf{1}[v_i = \text{scrutinize}]
\end{equation}
Under unconditional rejection without scrutiny, $\mathcal{C} = 0$ and
$c_i = 0$, so $u_i^{\text{reject}} = 0$.

\textit{Payoff under deviation.} If $a_i$ deviates to $v_i = 1$, and
$|\mathcal{B}| - 1 \geq |\mathcal{A}_t|/2$ still holds (which is true
when $|\mathcal{B}| > |\mathcal{A}_t|/2 + 1$), the outcome remains
rejection and the deviation yields no benefit. When $|\mathcal{B}|
= \lfloor |\mathcal{A}_t|/2 \rfloor + 1$, a deviation to $v_i = 1$
may flip the outcome to commitment, increasing $a_c$'s influence. Since
$u_i$ is strictly decreasing in $\mathbb{I}$:
\begin{equation}
    u_i^{\text{approve}} < u_i^{\text{reject}} = 0
\end{equation}
So deviation to approval is strictly dominated for all $a_i \in
\mathcal{B}$. Since $|\mathcal{B}|$ exceeds the blocking threshold,
$\mathcal{C}(\{v_i\}) = 0$ for all $\Delta_t$ from $a_c$, independently
of $\Delta_t$'s consistency with $\mathcal{M}_t$.
\end{proof}

\begin{remark}
\label{rem:failure_modes}
Propositions~\ref{prop:freeriding}--\ref{prop:adversarial_rejection}
are not mutually exclusive. A single auditor set $\mathcal{A}_t$ may
simultaneously contain free-riders (Proposition~\ref{prop:freeriding}),
a colluding sub-coalition (Proposition~\ref{prop:collusion}), and an
adversarially rejecting sub-coalition
(Proposition~\ref{prop:adversarial_rejection}). The three failure modes
operate on disjoint subsets of $\mathcal{A}_t$ and their effects
compound: free-riders provide no signal, colluders inject false approval,
and adversarial rejectors suppress truthful updates. Together, they
establish that no voting rule over untrusted auditors (e.g., majority,
supermajority, or weighted) can guarantee memory integrity in the
worst case, since each rule is vulnerable to at least one of the three
equilibria. This impossibility is the formal foundation for replacing
the vote-based commitment rule $\mathcal{C}$ with the structural
calibration mechanism $\Phi$ in Section~\ref{sec:calibration}.
\end{remark}



\textbf{Mapping to the three failure cases of Figure~\ref{fig:intro}.}
Propositions~\ref{prop:freeriding}--\ref{prop:adversarial_rejection} and the three failure cases described in Section~\ref{sec:intro} are related by a many-to-many correspondence rather than a one-to-one one, because each proposition characterizes an \emph{equilibrium structure} (which strategy profile is stable) while each case in Figure~\ref{fig:intro} characterizes an \emph{observed memory outcome} (what the resulting $\mathcal{M}_{t+1}$ looks like). A single equilibrium can produce different outcomes depending on the content of $\Delta_t$, and a single outcome can arise from different equilibria depending on which failure of agent confidence drives it. We make the correspondence explicit below.

\begin{table}[h]
\centering
\caption{Many-to-many correspondence between equilibrium failure modes (Propositions~\ref{prop:freeriding}--\ref{prop:adversarial_rejection}) and observed memory outcomes (Cases 1-3 of Figure~\ref{fig:intro}). \cmark{} indicates the proposition can produce the case; \xmark{} indicates it structurally cannot.}
\label{tab:case-prop-mapping}
\begin{tabular}{lccc}
\toprule
 & Case 1 & Case 2 & Case 3 \\
 & (corruption) & (over-rejection) & (uncertain/no commit) \\
\midrule
Prop.~\ref{prop:freeriding} (free-riding) & \cmark & & \cmark \\
Prop.~\ref{prop:collusion} (collusion) & \cmark & \cmark &  \\
Prop.~\ref{prop:adversarial_rejection} (adv.\ rejection) &  & \cmark & \cmark \\
\bottomrule
\end{tabular}
\end{table}

Proposition~\ref{prop:freeriding} (free-riding) produces Case 1 outcomes when the contributor is over-confident: rubber-stamp approval without scrutiny commits whatever $\Delta_t$ is proposed, including hallucinated entries with high stated confidence. It produces Case 3 outcomes when the contributor is itself hedged: rubber-stamp approval of a ``no-commit'' or noise proposal yields an under-populated memory. Free-riding does not produce Case 2, because rubber-stamping never rejects.

Proposition~\ref{prop:collusion} (collusion) produces Case 1 outcomes when colluding auditors approve a high-confidence corrupt contributor. This is the canonical ``two confident agents endorse a hallucinated entry'' pattern. It produces Case 2 outcomes when the colluding majority approves a confident contributor's \emph{rejection} of a tentative contribution, suppressing correct content with stated certainty. Collusion does not produce Case 3, because collusion is by definition active commitment, not abstention.

Proposition~\ref{prop:adversarial_rejection} (adversarial rejection) produces Case 2 outcomes directly: a blocking coalition vetoes truthful $\Delta_t$ with high confidence regardless of its alignment, which is exactly the over-curation pattern. It produces Case 3 outcomes when the blocking coalition's rejection coexists with weak contribution from the remaining agents (uncertain/no commit happens) and what could have been useful content is lost. Adversarial rejection does not produce Case 1, because rejection cannot inject corrupt content.


Overall, we have three observations: First, every case is reachable by at least two of the three equilibria, so no single proposition fully characterizes any single case. Second, every proposition produces at least two cases, so no single failure mode at the agent level produces a unique memory outcome. Third, the union of the three propositions covers all three cases, confirming that the failure-case taxonomy of Figure~\ref{fig:intro} is exhaustively grounded in the equilibrium analysis: any observable confidence-driven failure of shared-memory debate corresponds to at least one of the three propositions.

